\abstractPT  % Do NOT modify this line

Este trabalho documenta a implementa\c{c}\~{a}o do jogo de estrat\'{e}gia \textit{ISOLA} em
SWI-Prolog, no \^{a}mbito da unidade curricular de Intelig\^{e}ncia Artificial.
O jogo foi representado atrav\'{e}s de listas de listas em Prolog, com
predicados para gera\c{c}\~{a}o de movimentos, aplica\c{c}\~{a}o de jogadas e
detec\c{c}\~{a}o de fim de jogo. A intelig\^{e}ncia artificial do jogador
computorizado baseia-se no algoritmo \textit{alfa-beta}, uma implementa\c{c}\~{a}o
eficiente do princ\'{i}pio \textit{minimax}, complementada com ordena\c{c}\~{a}o
heur\'{i}stica de movimentos. A fun\c{c}\~{a}o de avalia\c{c}\~{a}o combina mobilidade
e centralidade. Foram realizados testes unit\'{a}rios (33 testes, todos bem
sucedidos) e benchmarks experimentais em tabuleiros de dimens\~{o}es e
profundidades variadas. Os resultados confirmam que a profundidade~4
\'{e} praticável em tabuleiros 3{\texttimes}3, mas exige tempo elevado
em tabuleiros 6{\texttimes}6 (cerca de 2~min~47~s no primeiro
movimento), justificando a profundidade~2 como compromisso para
tabuleiros maiores.
